\section{Health and safety impacts}
\label{sec:health_and_safety_impacts}

Health is also a critical dimension in \acp{bss} promotion, complementing the environmental and transport outcomes discussed earlier~\cite{DeMaio2009, Fishman2016, Qiu2018}. Regular cycling is consistently associated with improved cardiovascular fitness and lower risks of chronic disease and premature mortality~\cite{Green2021, Celis-Morales2017, Dutheil2020}, and it provides an efficient means of achieving the recommended 150 minutes of moderate to vigorous physical activity per week~\cite{Knowler2002}. Cycling can also support mental well-being~\cite{Chekroud2018, Ruiz2015}, but evidence that bike-sharing substantially increases physical activity at the population level is mixed. Early observational and survey-based studies reported increased cycling among residents near stations and self-reported health improvements among users~\cite{Fuller2013, Molina-Garcia2013, Alberts2012}, and county-level analyses linked \ac{bss} deployment with small declines in obesity prevalence~\cite{Xu2019, Hosford2019}. However, the magnitude of health gains depends strongly on which modes are replaced~\cite{Fishman2014, Murphy2015}. Accordingly, a narrative review concluded that many systems may have limited effects on overall population physical activity because \ac{bss} users are often already active and trips can substitute for walking or \ac{pt} rather than car travel~\cite{Bauman2017}.

To quantify net population-level effects and explicitly account for both benefits and risks, researchers commonly use \acp{hia}, which combine travel behavior data with epidemiological dose--response relationships~\cite{Rojas-Rueda2011, Woodcock2014, Otero2018, Clockston2021}. Here, \acp{hia} typically model three pathways~\cite{Teixeira2021}: (i) health gains from increased physical activity, (ii) health risks from traffic injuries, and (iii) health losses from exposure to air pollution, most often operationalized through in-transit \ce{PM_{2.5}} exposure~\cite{Woodcock2014, Otero2018, Clockston2021, Cao2021}. Results are sensitive to key assumptions, especially modal substitution, and to whether assessments account only for users' in-transit exposure or also include system-wide air-quality changes from reduced car use~\cite{Clockston2021, Hou2023}.

Across empirical applications, \acp{hia} consistently find that net health impacts are driven primarily by gains in physical activity, with air-pollution and traffic-injury pathways generally playing smaller roles. In Barcelona, Rojas-Rueda et al.~\cite{Rojas-Rueda2011} estimated that \textit{Bicing} prevented approximately 12.5 deaths per year, with comparatively small risks from traffic injuries and pollution exposure, yielding a benefit--risk ratio of about 77:1. However, this result partly reflected an optimistic modeling assumption of 90\% of \textit{Bicing} trips replacing car journeys, a limitation they acknowledged. In London, Woodcock et al.~\cite{Woodcock2014} similarly found a positive overall health balance using empirical trip data and disaggregation by age and sex, and reported heterogeneity across demographic groups: men experienced larger benefits because they cycled more frequently and for longer distances, older adults benefited more in absolute terms due to higher baseline mortality risk, and women faced higher relative injury risks in traffic.
 
Extending the evidence across 12 European systems, Otero et al.~\cite{Otero2018} reported a benefit--risk ratio of 19:1 and estimated 5.17 deaths avoided per 1,000 bicycles in service per year, while emphasizing that benefits could increase substantially under scenarios with higher car substitution. In the U.S., Clockston et al.~\cite{Clockston2021} estimated that \acp{bss} use prevented about 4.7 premature deaths annually nationwide (approximately \$36 million in avoided health costs) and around 2 premature deaths annually in New York City (about \$15 million). Beyond Europe and North America, Chen et al.~\cite{Chen2023} evaluated Shanghai’s \textit{Mobike} program in a highly polluted megacity context. For reference, in 2021, the \ac{who} recommended long-term average \ce{PM_{2.5}} concentrations not exceeding 5~\si{\micro g/m^3} and daily averages below 15~\si{\micro g/m^3}~\cite{WHO2021}. Despite average \ce{PM_{2.5}} concentrations of about 45~\si{\micro g/m^3}, their \ac{hia} found that health gains from increased physical activity still outweighed the added pollution burden, preventing roughly 7--8 premature deaths per year among about 306,000 users. The net balance began to diminish only as concentrations approached about 68~\si{\micro g/m^3}~\cite{Chen2023}.

The air-pollution pathway is typically marginal in standard \acp{hia} that focus on users' in-transit exposure~\cite{Rojas-Rueda2011, Woodcock2014, Otero2018, Clockston2021}, though it can become more influential in highly polluted contexts~\cite{Chen2023}. Some work has also emphasized that exposure can vary sharply over space and time, and that daily mean concentrations may obscure short-term peaks relevant to cyclists~\cite{Cao2021}.

Similarly, traffic-injury impacts generally remain small relative to physical-activity benefits, but they vary across cities and are sensitive to local safety environments and infrastructure~\cite{Rojas-Rueda2011, Woodcock2014, Otero2018, Clockston2021}. Evidence from London suggests that injury rates for bike-sharing users can be at least as low as those for private cyclists~\cite{Woodcock2014}. Surveys likewise report lower crash rates for bike-sharing than for private cycling and, in some analyses, lower risks of fatal or severe injuries after scheme introduction~\cite{Martin2016, Fishman2016_road_safety}. Indeed, reviews highlight a ``safety-in-numbers'' effect, whereby increases in cycling volumes can reduce the risk of injury per cyclist, helping to explain why higher ridership due to the existence of a \acp{bss} does not necessarily translate into proportionate increases in injuries~\cite{Jacobsen2015, Teixeira2021}. At the same time, a systematic review reported high rates of riding without a helmet among bike-sharing users and mixed evidence on whether injury rates change after scheme launches~\cite{Chen2021}. One multi-city analysis further found that in jurisdictions without mandatory helmet laws, the share of head-injured cyclists increased after bike-sharing implementation~\cite{Graves2014}. 

Finally, although \acp{bss} deliver clear population-level health benefits, these gains are not necessarily distributed equally across urban space or social groups. In New York City, station placement has historically favored higher-income neighborhoods, shaping who can access and benefit from the system; while network expansion increased estimated premature deaths averted (from two to three per year), sensitivity analyses suggested that prioritizing higher-poverty areas could have yielded larger benefits~\cite{Babagoli2019}. National-level \acp{hia} evidence from the U.S. similarly emphasizes that equity depends on more balanced station distribution and improved data on users’ socioeconomic characteristics~\cite{Clockston2021}. More broadly, because most evidence comes from high-income cities with relatively low air pollution and stronger road-safety environments, the benefit--risk balance may differ in settings with poorer air quality or weaker safety conditions, underscoring the need to expand the evidence base across contexts~\cite{Teixeira2021, Tainio2016}.